% Drafted Introduction and Related Work for paper on pFL bicuculline mapping.
% Used by writing-agent; inline into main.tex (do NOT \input{}).

\section{Introduction}

Breathing is a rhythmic motor behaviour generated by distributed neural oscillators in the ventral medulla. Inspiration depends on the preB{\"o}tzinger complex (preB{\"o}tC), a compact cluster of neurons whose genesis of rhythm has been characterised in detail over three decades \citep{Smith1991PreBotzinger,Feldman2013Understanding,delNegro2018Breathing}. Expiration, in contrast, is passive during eupnea but can become active when metabolic demand rises, airway resistance changes, or central chemosensory drive increases \citep{Janczewski2006Distinct,Iscoe1998Control}. Active expiration recruits the internal oblique and transverse abdominal muscles to drive airflow below the functional residual capacity, and its loss contributes to ventilatory insufficiency in conditions ranging from chronic obstructive pulmonary disease to motor neuron disease. A conditional expiratory oscillator located on the ventrolateral surface of the medulla, rostral to the preB{\"o}tC and overlapping the caudal aspect of the facial nucleus, is now widely accepted as the primary substrate of this active phase \citep{Pagliardini2011Active,Pisanski2019Parafacial,delNegro2018Breathing}.

This region is referred to as the lateral parafacial nucleus (pFL) and lies ventral and lateral to the retrotrapezoid nucleus (RTN), the well-studied Phox2b-expressing chemosensory nucleus that modulates breathing in response to central CO\textsubscript{2} and pH \citep{Guyenet2010Central,GuyenetAbbott2013RTN,Takakura2013Phox2b,Marina2010Essential}. Pharmacological disinhibition of the pFL with GABA-A receptor antagonists, in particular with bicuculline or the mixed GABA-A/glycine antagonist bicuculline/strychnine, reliably evokes active expiration in anaesthetized adult rats and in situ preparations, accompanied by phasic late-expiratory (late-E) bursts in the integrated abdominal electromyogram \citep{Pagliardini2011Active,Abdala2009Abdominal,Molkov2010Computational,deBritto2017GABAergic}. Chemogenetic inhibition of pFL neurons, conversely, reduces the intensity of expiratory muscle recruitment and diminishes but does not abolish active expiration during sleep in rats \citep{Huckstepp2015Role,Pisanski2020Parafacial}. The pFL therefore behaves as a conditional oscillator: permissive for active expiration, but not obligatory for eupneic breathing.

Despite broad agreement on the functional identity of pFL, its precise anatomical core has remained surprisingly debated. Different groups have targeted this region at rostrocaudal coordinates ranging from -0.3 to +1.0 mm relative to the caudal tip of the facial nucleus (VIIc), the common anatomical landmark on the ventral medulla \citep{Pagliardini2011Active,Huckstepp2015Role,deBritto2017GABAergic,Silva2019Importance,Boutin2017Respiratory}. In juvenile rats, late-expiratory activity and ABD/DIA coupling have been reported after bicuculline injections at +0.1 to +0.3 mm from VIIc \citep{deBritto2017GABAergic}. In adult rats, chemogenetic and pharmacological targeting at +0.5 mm produces robust ABD recruitment \citep{Huckstepp2015Role,Pagliardini2011Active}. A separate line of work, using bicuculline/strychnine in combination with hypercapnia, localises the core more rostrally, at +0.3 to +1.0 mm from VIIc \citep{Silva2019Importance}. The qualitative pattern of the ABD response also varies across these studies: purely tonic expiratory bursts have been reported with combined GABAergic and glycinergic disinhibition under CO\textsubscript{2}, whereas bicuculline alone in normocapnia tends to produce a mixture of tonic and late-E patterns \citep{Pagliardini2011Active,deBritto2017GABAergic,Silva2019Importance}. Whether these differences reflect species, age, coordinate, anaesthetic, gas composition, or a genuine microzonal architecture of pFL has been impossible to adjudicate from a meta-analysis of these heterogeneous studies.

A further question is how active expiration interacts with the rest of the respiratory cycle when it is driven at different rostrocaudal coordinates. Active expiration not only recruits abdominal muscles but also modifies inspiratory timing, inspiratory airflow, post-inspiratory glottal control, and global minute ventilation through effects on diaphragm recruitment, upper airway resistance, and sympathetic drive \citep{Molkov2010Computational,Dutschmann2006Pontine,Zoccal2018Expiratory,Barnett2017Partial,Dhingra2020LateE}. Most prior studies analyse the ABD response with univariate summaries (peak amplitude, duration, frequency), which collapse the multi-phase structure of the respiratory cycle. A multivariate analysis that combines airflow and the integrated DIA and ABD EMG signals can separate contributions from late-E, inspiratory, and post-inspiratory windows of the cycle, and may therefore reveal microzonal differences along the rostrocaudal axis of pFL.

Here, we provide a systematic within-study functional map of pFL in the anaesthetized, vagotomized adult rat. We bilaterally pressure-injected bicuculline (200 $\mu$M, 200 nL per side) at one of five coordinates spanning -0.2 mm caudal to +0.8 mm rostral of VIIc, or HEPES vehicle, and recorded ABD and DIA EMG, airflow, tidal volume, minute ventilation, and oxygen consumption for twenty minutes. We verified each injection site using fluorescent beads and cFos/PHOX2B immunohistochemistry \citep{Marina2010Essential,Takakura2013Phox2b}. We then quantified the response with univariate time series and with a multivariate 3D trajectory of the respiratory cycle. The results show that active expiration is elicitable at every coordinate tested, but that small shifts along the rostrocaudal axis ($\leq$400 $\mu$m) qualitatively reorganise the pattern: the fastest, longest, and most metabolically consequential responses arise at +0.6 to +0.8 mm from VIIc, and the +0.6 and +0.8 mm microzones themselves differ in which phase of the cycle they dominate. These findings reconcile the heterogeneity of prior coordinates by grounding the expiratory oscillator in a compact rostral core, and they sharpen the coordinates used by subsequent physiological and translational studies of active expiration \citep{Pagliardini2024Parafacial,Jenkin2017Roles,Pisanski2020Parafacial,Pisanski2019Parafacial}.
